\section{感知机：深度学习的基本构建块}

\subsection{感知机的前向传播}

\textbf{感知机}（Perceptron）是所有神经网络的基本结构单元，也被称为一个"神经元"（neuron）。其前向传播过程包含三个步骤：

\begin{enumerate}
\item \textbf{加权求和}：将每个输入 $x_i$ 乘以对应权重 $w_i$，然后求和
\item \textbf{加偏置}：加上偏置项 $w_0$（bias）
\item \textbf{非线性激活}：通过非线性激活函数 $g(\cdot)$ 得到输出
\end{enumerate}

\begin{figure}[H]
\centering
\includegraphics[width=0.85\textwidth]{slides-images/slide-022.jpg}
\caption{感知机的前向传播（含偏置项）\protect\footnotemark}
\end{figure}
\footnotetext{Slides 第 22 页。}

其数学表达式为：

$$\hat{y} = g\left(w_0 + \sum_{i=1}^{m} x_i w_i\right) = g(w_0 + \mathbf{X}^T \mathbf{W})$$

其中：
\begin{itemize}
\item $\mathbf{X} = [x_1, x_2, \ldots, x_m]^T$：输入向量
\item $\mathbf{W} = [w_1, w_2, \ldots, w_m]^T$：权重向量
\item $w_0$：偏置项（bias），允许激活函数沿 $z$ 轴平移
\item $g(\cdot)$：非线性激活函数
\item $\hat{y}$：预测输出
\end{itemize}

\begin{importantbox}{偏置项的作用}
偏置项 $w_0$ 是感知机中不可或缺的参数。它允许激活函数独立于输入值进行平移，使得即使所有输入为零，神经元也能产生非零输出。没有偏置的感知机只能学习通过原点的决策边界。
\end{importantbox}

\subsection{激活函数}

激活函数 $g(\cdot)$ 的核心作用是\textbf{为网络引入非线性}。常见的激活函数包括：

\begin{figure}[H]
\centering
\includegraphics[width=0.95\textwidth]{slides-images/slide-025.jpg}
\caption{三种常见激活函数：Sigmoid、Tanh 和 ReLU\protect\footnotemark}
\end{figure}
\footnotetext{Slides 第 25 页。}

\begin{table}[H]
\centering
\begin{tabular}{llll}
\toprule
\textbf{激活函数} & \textbf{公式} & \textbf{输出范围} & \textbf{特点} \\
\midrule
Sigmoid & $\sigma(z) = \frac{1}{1+e^{-z}}$ & $(0, 1)$ & 适合概率输出 \\
Tanh & $g(z) = \frac{e^z - e^{-z}}{e^z + e^{-z}}$ & $(-1, 1)$ & 零中心化 \\
ReLU & $g(z) = \max(0, z)$ & ${[}0, +\infty)$ & 计算高效，最常用 \\
\bottomrule
\end{tabular}
\caption{常见激活函数对比}
\end{table}

\begin{warningbox}{为什么必须使用非线性激活函数？}
如果不使用非线性激活函数（即使用线性激活 $g(z)=z$），那么无论网络有多少层，整个网络的输出仍然只是输入的线性组合。这意味着多层网络退化为单层线性模型，无法学习任何非线性模式。只有引入非线性，网络才能逼近任意复杂的函数。
\end{warningbox}

\begin{figure}[H]
\centering
\includegraphics[width=0.85\textwidth]{slides-images/slide-028.jpg}
\caption{非线性激活函数使神经网络能够学习复杂的非线性决策边界\protect\footnotemark}
\end{figure}
\footnotetext{Slides 第 28 页。左：线性激活只能产生线性决策边界；右：非线性激活可以逼近任意复杂的函数。}

\subsection{感知机的几何直觉}

Amini 通过一个具体的二维示例来建立直觉。考虑一个具有两个输入 $(x_1, x_2)$ 的感知机，参数为 $w_0=1$，$\mathbf{W}=[3, -2]^T$：

$$\hat{y} = g(1 + 3x_1 - 2x_2)$$

\begin{figure}[H]
\centering
\includegraphics[width=0.85\textwidth]{slides-images/slide-032.jpg}
\caption{感知机示例：决策边界将二维空间分为两个区域\protect\footnotemark}
\end{figure}
\footnotetext{Slides 第 32 页。}

括号内的线性部分 $z = 1 + 3x_1 - 2x_2$ 定义了二维空间中的一条直线。这条线就是\textbf{决策边界}：
\begin{itemize}
\item 当 $z > 0$ 时，经过 Sigmoid 激活后 $\hat{y} > 0.5$，预测为正类
\item 当 $z < 0$ 时，$\hat{y} < 0.5$，预测为负类
\item 当 $z = 0$ 时，$\hat{y} = 0.5$，恰好在决策边界上
\end{itemize}

以输入 $\mathbf{X}=[-1, 2]^T$ 为例：$\hat{y} = g(1 + 3\times(-1) - 2\times 2) = g(-6) \approx 0.002$，远小于 0.5，因此该点被分类为负类。

\begin{knowledgebox}{单个感知机的局限}
单个感知机只能学习线性决策边界（在二维空间中是一条直线，在高维空间中是一个超平面）。这意味着它无法解决 XOR 等非线性可分问题。要处理更复杂的模式，需要将多个感知机组合成多层网络。
\end{knowledgebox}

\subsection{本章小结}

感知机是神经网络的最小单元，其前向传播包括加权求和、加偏置、通过激活函数三步。偏置项允许决策边界不必通过原点，非线性激活函数使网络能学习复杂模式。单个感知机只能表示线性分类器。

